\section{Compound Exposure Corridors}
\label{sec:compound-exposure}

\subsection{Identification Methodology}

The compound resilience hole condition requires simultaneously elevated scores on both the B1 (isolation) and D1 (climate vulnerability) dimensions of the ROUTE rubric \citep{ROUTE_A2}. We use thresholds of B1 $>$ 7 AND D1 $>$ 6, which correspond to the 85th percentile of each dimension in the v1.1 corpus of 227 T1 and T2 corridors. The B1 score is defined as the inverse of alternate-route adequacy: a corridor scores B1 = 10 if closure produces national-scale freight disruption with no interstate-standard alternate. The D1 score is defined as composite climate vulnerability: probability-weighted exposure to events causing multi-hour closures, including winter precipitation events, hurricane-force wind and surge, riverine flooding, and wildfire smoke or evacuation closures.

Applying both thresholds to the v1.1 corpus identifies 11 corridors in compound exposure. Of these, 4 score B1 $>$ 8 AND D1 $>$ 7 (extreme compound exposure) and are the primary focus of this section. Table~\ref{tab:compound} presents the full compound exposure matrix.

\begin{table}[h!]
\centering
\caption{Compound Resilience Hole Corridors: B1 and D1 Scores}
\label{tab:compound}
\begin{tabular}{lrrrr}
\toprule
\textbf{Corridor} & \textbf{B1} & \textbf{D1} & \textbf{Closure events/yr} & \textbf{Annual cost (\$B)} \\
\midrule
I-80 Donner Pass (CA)        & 8.3 & 7.8 & $\sim$50 & 1.60 \\
I-90 Snoqualmie Pass (WA)    & 7.6 & 7.1 & $\sim$40 & 0.65 \\
I-10 Gulf Coast LA           & 6.8 & 8.4 & $\sim$8  & 0.82 \\
I-10 Gulf Coast TX (Houston) & 5.9 & 7.6 & $\sim$6  & 0.54 \\
I-35 Oklahoma Tornado Corr.  & 6.2 & 7.2 & $\sim$12 & 0.31 \\
I-5 Siskiyou Pass (OR)       & 7.1 & 6.8 & $\sim$30 & 0.48 \\
I-90 Homestake Pass (MT)     & 7.8 & 6.4 & $\sim$25 & 0.22 \\
\bottomrule
\end{tabular}
\end{table}

\subsection{I-80 Donner Pass: Extreme Compound Exposure}

Donner Pass is the canonical compound resilience hole. Interstate 80 at the Sierra Nevada crossing (elevation 7,227 feet) scores B1 = 8.3, reflecting the absence of any interstate-standard alternate across the Sierra Nevada between I-40 (central California via Barstow, 310 miles south) and the Oregon border. The D1 score of 7.8 reflects an average of 50 closure events per year, with a mean duration of 18 hours and maximum duration exceeding 72 hours \citep{Caltrans2023}.

The combination produces a freight disruption exposure that is unmatched in the national network. Approximately 8,000 commercial vehicles per day traverse Donner Pass during peak freight season. When the pass closes, these vehicles face a binary choice: wait (losing 18 hours on average) or reroute (adding 310 miles and 5.5 hours of driving, plus a \$225/hr premium for fuel, driver, and schedule penalty \citep{ATRI_costs2024}). Because neither option is economically acceptable for just-in-time freight, the compound effect is not merely additive --- it multiplies. A shipper managing a 24-hour delivery window cannot absorb either a mean 18-hour wait or a 5.5-hour detour without a service failure.

Sierra Nevada snowpack is projected to decrease in total depth but increase in storm intensity variability under RCP 8.5 climate scenarios \citep{NOAA_ClimateAtlas2022}, meaning fewer moderate-accumulation events but more extreme precipitation events. The D1 trajectory for Donner Pass is therefore toward higher-severity, lower-frequency closures --- each closure longer in duration --- rather than reduced exposure. The 2024 D1 score of 7.8 is a baseline that the ROUTE climate model projects rising to 8.5 by 2050.

\subsection{I-10 Gulf Coast Louisiana: High D1 Compound Exposure}

The I-10 corridor through southern Louisiana presents a different compound exposure profile. The B1 score is moderate (6.8) rather than extreme: I-10 has a functional alternate via I-20 through northern Louisiana, but the alternate adds 180 miles and does not serve the Port of New Orleans or Baton Rouge industrial complex, which generates approximately 12\% of US petrochemical freight. The D1 score is the highest in the national corpus for a major T1 corridor (8.4), reflecting the combination of: (1) mean sea level 1.2--1.8 meters below regional hurricane surge levels, (2) riverine flooding from the Mississippi and Atchafalaya systems during high-water years, and (3) subsidence-driven pavement degradation that accelerates closure risk.

The Gulf Coast I-10 compound exposure is distinct from Donner in its climate trajectory. Donner's risk is wintertime precipitation. Gulf Coast I-10's dominant risk by 2050 is sea level rise and hurricane surge intensity. NOAA projects 0.5--1.0 meter of additional sea level rise in the Louisiana Gulf Coast by 2050 under intermediate scenarios \citep{NOAA_SLR2022}. This would raise the annual probability of a major surge inundation event from approximately 8 per year (Category 1--2 direct hits) to 12--15 per year, increasing the D1 score from 8.4 to a projected 9.1 by 2050.

The I-10 Gulf Coast investment case is different from Donner. A freight tunnel is not economically viable or geologically appropriate across Louisiana bayou country. The intervention is infrastructure hardening: roadway elevation at the six lowest-lying segments (mean elevation 0.8m above sea level), improved drainage structures, and hurricane-resistant barrier design. Estimated cost: \$2.1 billion for the Louisiana segment. The B1 component is addressed through I-20 and I-10B designated freight alternate corridors with adequate load rating and managed interchange access.

\subsection{I-90 Snoqualmie Pass: Sierra-Analog Compound Exposure}

Interstate 90 at Snoqualmie Pass (elevation 3,022 feet) is structurally analogous to Donner but at lower severity. B1 = 7.6 reflects the absence of a viable commercial truck alternate across the Washington Cascades during winter closures. US-2 over Stevens Pass provides 60 additional miles but is not maintained for commercial freight during closures \citep{WSDOT2023}. D1 = 7.1 reflects approximately 40 annual closure events with mean duration of 12 hours.

The Snoqualmie compound exposure is particularly consequential for Pacific Northwest agricultural freight: the Yakima Valley produces approximately 75\% of US commercial apple output, and I-90 is the primary route from the valley to Seattle ports. Closure events during October--November harvest windows impose both the direct rerouting cost and a commodity spoilage cost not captured in standard freight delay metrics.

\subsection{I-35 Oklahoma Tornado Corridor: Low-B1, Moderate-D1 Compound Exposure}

The I-35 corridor through central Oklahoma presents a compound exposure profile driven primarily by tornado and severe convective weather events rather than winter precipitation or flooding. B1 = 6.2 is moderate: US-77 and US-81 provide parallel alignment but are not adequate for commercial freight volumes. D1 = 7.2 reflects approximately 12 significant weather closure events per year, including tornado warning closures (mandatory evacuation of interstate for shelter-in-place protocol), high-wind closures affecting high-profile vehicles, and hail events that close segments for emergency response.

The investment case for I-35 compound exposure is the lowest of the priority corridors --- partial alternate route upgrade (US-77 to interstate standard between Oklahoma City and Wichita, approximately \$1.8B) rather than a primary corridor hardening --- but the compound exposure score of 6.2 $\times$ 7.2 = 44.6 (using the B1$\times$D1 product metric) places it within the priority investment zone.

\subsection{Investment Sequencing from Compound Exposure}

The compound exposure matrix produces a different investment priority order than either B1 or D1 alone. Sorted by the B1$\times$D1 product score, the priority sequence is: Donner Pass (64.7), Gulf Coast I-10 Louisiana (57.1), Snoqualmie (54.0), Siskiyou (48.3), I-35 Oklahoma (44.6). This sequence differs from the D1-only ordering (which would prioritize Gulf Coast I-10 above Donner) and from the B1-only ordering (which would prioritize Donner and Snoqualmie but underweight Gulf Coast). The compound metric identifies the corridors where both structural isolation and environmental exposure are simultaneously acute --- and where investment in either dimension alone is insufficient.
